\chapter{Conclusion}
\label{chap:conclusion}

\section{Summary}
We studied training-free Video Anomaly Detection on industrial conveyor scenes with
Vision-Language Models through three approaches. Direct prompting established baselines
and exposed the model's bias toward ``normal''. Multi-persona majority voting on a local
open-weight model addressed this bias through prompt diversity and the articulate-first
format. VERA-lite optimised the guiding questions, introduced a blind evaluation
protocol, compared input representations, and combined two views into an ensemble that
reached $13/13$ on the pen dataset with zero false positives.

\section{Key Findings}
\begin{itemize}
  \item \textbf{Articulation / question decomposition counters the ``normal''-bias.}
        Forcing the model to describe what it sees before deciding --- whether via the
        articulate-first persona format or via VERA guiding questions --- improves
        reliability.
  \item \textbf{A continuous-vs-discrete trade-off is intrinsic for temporal anomalies.}
        Continuous video captures edge-orientation and belt-stop cues but misses subtle
        intra-object motion; discrete frame representations capture the motion but lose
        the other cues. The trade-off held across three discrete representations.
  \item \textbf{Ensembling recovers the union.} Because the video and dense-frame views
        have disjoint blind spots and zero false positives, their OR-ensemble reaches
        $13/13$ at no extra precision cost.
  \item \textbf{Single-run accuracy is misleading.} Repeated runs put the true accuracy
        at $0.846$ ($11/13$ in all three runs), not the lucky single-run $0.92$; the
        temporal limit is deterministic while two borderline clips flip between runs
        under the unbounded baseline but not under a capped reasoning budget.
  \item \textbf{Reasoning-budget control is a safe speed lever.} Capping the reasoning
        budget gave a $\sim$41\% speed-up at identical accuracy and made the per-clip
        verdicts deterministic, whereas raising the sampling rate was counter-productive.
\end{itemize}

\section{Limitations}
The pen-conveyor dataset is a $13$-clip pilot recorded under controlled, fixed
conditions; reaching $13/13$ by tuning on these exact clips risks overfitting, and the
controlled setup hides robustness to lighting, angle and resolution changes. Only one
API-served model was reliably accessible. The three approaches used different models and
scenes, so their numbers are not directly comparable. The persona approach lacks a
multi-clip quantitative evaluation.

\section{Future Work}
Priorities are: (i) a larger dataset with a proper train/test split and statistical
confidence intervals; (ii) simple computer-vision baselines (frame differencing, optical
flow) to quantify what the VLM adds; (iii) frame-level metrics (AUROC/AP) using the
available frame-level ground truth; (iv) a fair head-to-head evaluation of all three
approaches on the same labelled set, including the three persona ablations; (v) a
calibration analysis, since the model can be confidently wrong; and (vi) a cascade
architecture (cheap detector first, VLM only on flagged segments) for real-time
deployment.
